%% kapa: exam=1 round=22 year=2011 period=1 form=A subject=civil_law subject_ko=민법 seq=34 number=34
%% source: Q-Net 감정평가사 기출문제 (cst00309) artlSeq=5036790 file=extracted/1008/5036790/2011년 제22회 감정평가사 1차시험 기출문제/1교시.hwp
%% conversion: kordoc 4.24.51 -> tools/gen1.py (auto); review=proofread-sonnet; flags=none
\begin{question}{34}
등기에 관한 설명으로 옳지 않은 것은? (다툼이 있으면 판례에 의함)
\begin{choices}
  \item 하나의 부동산에 관하여 동일인 명의로 중복보존등기가 경료된 경우, 뒤에 경료된 등기는 무효이다.
  \item 등기가 멸실된 경우 소정기간 내에 멸실회복등기가 없는 경우에도 소유권은 소멸하지 않는다.
  \item 법률행위로 인한 등기청구권은 원칙적으로 소멸시효에 걸리지 않는다.
  \item 경정등기는 등기와 실체관계 사이의 원시적 불일치를 시정하기 위해서 하는 등기를 말한다.
  \item 등기명의인 표시변경등기는 등기명의인의 동일성이 유지되는 범위 내에서 등기부상의 표시를 실제와 합치시키기 위하여 행하여지는 것에 불과할 뿐 어떠한 권리변동을 가져오는 것은 아니다.
\end{choices}
\end{question}
\answer{3}
